% ECONOMICS_PROBLEM: {"id": "C73-27", "title": "Large-population selection of finite-state potential minimizers", "jel_code": "C73", "source_id": "OP-0214"}
% !TeX program = xelatex
\documentclass[11pt,a4paper]{article}
\usepackage[margin=23mm]{geometry}
\usepackage{fontspec}
\setmainfont{Latin Modern Roman}
\usepackage{amsmath,amssymb,mathtools}
\usepackage{xcolor,enumitem,booktabs,longtable,array,xurl,hyperref}
\definecolor{muted}{HTML}{53636C}
\hypersetup{colorlinks=true,urlcolor=blue,linkcolor=blue}
\setlength{\parindent}{0pt}
\setlength{\parskip}{5pt}
\newcommand{\R}{\mathbb R}
\newcommand{\E}{\mathbb E}
\newcommand{\N}{\mathbb N}
\newcommand{\ind}{\mathbf 1}
\newcommand{\Var}{\operatorname{Var}}
\newcommand{\Cov}{\operatorname{Cov}}
\newcommand{\supp}{\operatorname{supp}}
\DeclareMathOperator*{\argmin}{arg\,min}
\DeclareMathOperator*{\argmax}{arg\,max}
\newcommand{\entrymeta}[3]{{\sffamily\small\color{muted}\textbf{\texttt{#1}}\quad|\quad #2\par #3\par}}
\newcommand{\Problemref}[1]{\texttt{#1}}
\newcommand{\JELref}[2]{\texttt{#2} (JEL \texttt{#1})}
\begin{document}
\section*{Large-population selection of finite-state potential minimizers}
\textbf{Problem ID:} \texttt{C73-27}\quad \textbf{JEL:} \texttt{C73}\par
% BEGIN ECONOMICS STATEMENT
\label{problem:OP-0214}
% GAME_THEORY_PROBLEM: {"local_id": "GT-084", "original_number": 84, "kind": "P", "entry_kind": "statement_only", "published": false, "review_required": true, "category": "game_theory"}
\label{game:GT-084}
{\sffamily\small\color{muted}
{\small\textbf{GT-084} | Mean-field games | P: published selection question; source-date status only\par}
\par}
\subsection*{Definitions and mathematical statement}
Let $E=\{1,\ldots,d\}$, $d\geq3$, $\Delta_d=\{m\in[0,1]^d:\sum_i m_i=1\}$, $T>0$, and $m_0\in\operatorname{int}\Delta_d$. Fix $C^3$ potentials $\mathcal F,\mathcal G$ on a neighborhood of $\Delta_d$, and set
\[
 f_i(m)=\partial_{m_i}\mathcal F(m),\qquad
 g_i(m)=\partial_{m_i}\mathcal G(m).
\]
Using extensions merely specifies component representatives: tangent differences of these gradients determine the rates on the simplex.

An admissible deterministic control is a measurable rate matrix $a(t)$ with $a_{ij}\geq0$ for $j\ne i$ and $a_{ii}=-\sum_{j\ne i}a_{ij}$, together with an absolutely continuous law $m:[0,T]\to\Delta_d$ satisfying
\[
 \dot m_i=\sum_{j\ne i}(m_j a_{ji}-m_i a_{ij}),\qquad m(0)=m_0,
 \qquad
 \int_0^T\sum_{i}\sum_{j\ne i}m_i a_{ij}^2dt<\infty.
\]
Its potential action is
\[
 \mathcal J(m,a)=\int_0^T\left[
 \tfrac12\sum_i\sum_{j\ne i}m_i(t)a_{ij}(t)^2+\mathcal F(m(t))\right]dt
 +\mathcal G(m(T)).
\]
Let $\mathcal M_{T,m_0}$ be the set of law paths $m$ occurring in global minimizers of this action, with the same initial condition. This is a cooperative potential functional, not $\sum_i m_if_i$ in place of $\mathcal F$.

In the associated $N$-player game, player $\ell$ has state $X^\ell\in E$ and controls its jump rates from its current state $i$ to $j\ne i$. Let $m_t^{N,-\ell}=(N-1)^{-1}\sum_{r\ne\ell}\delta_{X_t^r}$. Its cost is
\[
 \mathbb E\!\left[\int_0^T\left(
 \tfrac12\sum_{j\ne X_t^\ell}|a_{X_t^\ell j}^{\ell}(t)|^2
 +f_{X_t^\ell}(m_t^{N,-\ell})\right)dt
 +g_{X_T^\ell}(m_T^{N,-\ell})\right].
\]
Admissible finite-player rates are bounded progressively measurable matrices with nonnegative off-diagonal entries and row sums zero. Individual transition noises are independent, and there is no common noise. Consider bounded symmetric Markov-perfect Nash feedbacks: at every time and state vector, no player can reduce its continuation cost by changing its own admissible rates while the other feedback functions stay fixed. Initial states are independent with common law $m_0$ and independent of future transition noises. Write $\mu_t^N=N^{-1}\sum_\ell\delta_{X_t^\ell}$.

\textbf{Question.} If along any sequence of these equilibria
\[
 \mathcal L(\mu^N)\Longrightarrow Q
 \quad\text{in }D([0,T],\Delta_d)
 \quad\text{and}\quad Q(C([0,T],\Delta_d))=1,
\]
must $Q(\mathcal M_{T,m_0})=1$? The topology on $D$ is Skorokhod $J_1$. A deterministic limit is the special case $Q=\delta_m$. The question concerns globally minimizing paths, rather than merely first-order equilibria.
\subsection*{Short English statement}
Must every equilibrium selected by the large finite game minimize the corresponding potential control problem when there are at least three states?
\subsection*{Variants and scope boundaries}
The deterministic MFG conditions are
\[
 -\dot u_i+\tfrac12\sum_{j\ne i}(u_i-u_j)_+^2=f_i(m),\quad
 \dot m_i=\sum_{j\ne i}[m_j(u_j-u_i)_+-m_i(u_i-u_j)_+],\quad
 u_i(T)=g_i(m(T)).
\]
Here $r_+=\max(r,0)$. These are necessary optimality/equilibrium conditions, not a global-minimum test. The source reports a positive two-state selection result; finite-player common noise, open-loop equilibria, and arbitrary approximate equilibria are separate variants.
\subsection*{Source relationship and status}
Section~4.3 explicitly separates the open $d\geq3$ large-$N$ question from its two-state case. The path-law formulation allows randomized subsequential selection. No exhaustive later-literature audit is claimed.
\subsection*{References}
\begingroup\footnotesize\raggedright
{\footnotesize\raggedright
\textbf{[S1]} Cardaliaguet and Delarue. \emph{Selected topics in mean field games}, ICM2022, vol.~5; DOI: 10.4171/ICM2022/17.\par
\textit{Locators:} (4.7)--(4.10), pp.~3683--3684, rate and equilibrium definitions; (4.18), p.~3689, potential condition; (4.19), Proposition~4.4 and ``Selection of equilibria,'' p.~3690, action and explicit $|E|\geq3$ question.\par
\url{https://ems.press/content/book-chapter-files/33265}\par}
\par\endgroup

\subsection*{Common conventions: Source mathematical conventions}

\textbf{Finite games.} For a finite set $A$, $\Delta(A)$ is its probability
simplex. Players choose independent mixed strategies in Nash equilibrium;
correlation is allowed only when explicitly part of the solution concept.
In a game with payoff functions $u_i$ and strategy profile $x$, an additive
$\varepsilon$ Nash equilibrium satisfies
\[
 u_i(x)\geq u_i(x_i',x_{-i})-\varepsilon
 \quad\text{for every player $i$ and every admissible deviation $x_i'$}.
\]
For numerical approximation, payoffs are normalized as locally specified.
An $\varepsilon$ payoff residual is distinct from distance at most
$\varepsilon$ to the set of exact equilibria. Point convergence, convergence
of empirical play, and convergence in distance to an equilibrium set are
also distinct.

\textbf{Computational representations.} Unless stated otherwise, a complexity
question takes rational data with binary encoding; polynomial time is measured
in total bit length. A fixed-accuracy polynomial algorithm is not necessarily
polynomial in $1/\varepsilon$, and neither is necessarily polynomial in
$\log(1/\varepsilon)$. Succinct games and value, demand, or sampling oracles
must declare their interfaces. Exact real solutions require an explicit
representation; an unspecified list of arbitrary real numbers is not a
finite computational output. A claimed algorithmic barrier is meaningful
only for its specified model and reductions.

\textbf{Dynamic games.} Histories, information, observation, and allowable
randomization are part of the model. A stationary strategy depends only on
the current state; a positional strategy in a deterministic graph game
chooses one action at each controlled vertex. Nash equilibrium at an initial
state and equilibrium after every history are different requirements.
An expectation of a pathwise $\liminf$ payoff, a $\liminf$ of expected averages,
a limiting expected average, and a uniform finite-horizon equilibrium are
different criteria. None is silently substituted for another.

\textbf{Allocation and mechanisms.} A complete allocation assigns every item
exactly once unless otherwise stated. Zero-valued goods, negative values,
capacity constraints, feasibility, lotteries, and copies of goods affect
fairness definitions and are specified locally. Dominant-strategy incentive
compatibility, Bayesian incentive compatibility, universal truthfulness,
and truthfulness in expectation impose different inequalities. Whenever
an objective ratio has zero denominator, its convention is stated or those
instances are excluded.

\textbf{Infinite-dimensional models.} Expectations, integrals, stochastic
equations, and optimization objectives require their local regularity,
integrability, filtrations, and admissible domains. A backward--forward
mean-field system is a boundary-value problem; it is not an initial-value
semiflow without an additional construction. Long-time, large-population,
and vanishing-noise limits require an order of limits and a convergence
topology. If a minimum need not be attained, the objective is its infimum
or supremum and the approximation requirement is stated separately.
% END ECONOMICS STATEMENT
\end{document}
